\uuid{TESq}
\titre{Probabilités à un, deux et trois écarts-types}
\auteur{Maxime Nguyen}
\organisation{AMSCC}
\datecreate{2026-11-27}
\niveau{L2}
\module{Probabilité}
\chapitre{Probabilité continue}
\sousChapitre{Loi normale}
\difficulte{1}
\theme{loi normale, standardisation, ordres de grandeur}

\contenu{
\begin{question}
On note $\mathcal{N}(\mu,\sigma)$ la loi normale d'espérance $\mu$ et d'écart-type $\sigma$, et $\Phi$ la fonction de répartition de $\mathcal{N}(0,1)$. On donne $\Phi(1) \approx 0{,}841$, $\Phi(2) \approx 0{,}977$ et $\Phi(3) \approx 0{,}9987$. Soit $X \sim \mathcal{N}(100, 15)$. Quelles affirmations sont vraies ?
\begin{answers}
\good{$\mathbb{P}(85 \leq X \leq 115) \approx 0{,}68$.}
\bad{$\mathbb{P}(X \geq 130) \approx 0{,}05$.}
\good{$\mathbb{P}(70 \leq X \leq 130) \approx 0{,}95$.}
\good{$\mathbb{P}(X \leq 55) \approx 0{,}0013$.}
\bad{$\mathbb{P}(X \leq 115) \approx 0{,}68$.}
\end{answers}
\begin{explanations}
Avec $Z = \frac{X - 100}{15}$ de loi $\mathcal{N}(0,1)$ :

A : $\mathbb{P}(-1 \leq Z \leq 1) = 2\Phi(1) - 1 \approx 0{,}68$.

B : $\mathbb{P}(Z \geq 2) = 1 - \Phi(2) \approx 0{,}023$. La valeur $0{,}05$ est la probabilité de sortir de $[70;130]$ des deux côtés.

C : $\mathbb{P}(-2 \leq Z \leq 2) = 2\Phi(2) - 1 \approx 0{,}95$.

D : $\mathbb{P}(Z \leq -3) = 1 - \Phi(3) \approx 0{,}0013$.

E : $\mathbb{P}(Z \leq 1) = \Phi(1) \approx 0{,}84$ ; $0{,}68$ est la probabilité de l'intervalle centré.

À retenir : environ $68\,\%$, $95\,\%$ et $99{,}7\,\%$ de la masse à moins de
un, deux et trois écarts-types de l'espérance.

Contre-sens visés : confondre une queue et les deux queues (B) ; confondre
$\Phi(t)$ et $2\Phi(t) - 1$ (E).
\end{explanations}
\end{question}
}
